% Figure: skin depth of copper versus frequency, log-log Cartesian axes.
% delta = sqrt(2 / (omega mu sigma)) = 1/sqrt(pi f mu sigma).
% For copper sigma = 5.96e7 S/m, mu = mu0 = 4 pi e-7, so
% delta(f) = 0.0660 / sqrt(f) metres (f in Hz), i.e. 66.0 / sqrt(f) mm.
% We plot log10(delta in micrometres) versus log10(f in Hz) as a straight line.
% delta_um = 6.60e7 / sqrt(f) micrometres; log10(delta_um) = 7.82 - 0.5 log10(f).
\begin{figure}[htbp]
\centering
\begin{tikzpicture}
\begin{axis}[
  width=12cm, height=6cm,
  xlabel={$\log_{10}$ frequency (Hz)},
  ylabel={$\log_{10}$ skin depth ($\mu$m)},
  xmin=3, xmax=11, ymin=0, ymax=5,
  domain=3:11, samples=2,
  axis lines=left,
  xtick={3,4,5,6,7,8,9,10,11},
  ytick={0,1,2,3,4,5},
  enlargelimits=false,
]
  % straight line: y = 7.82 - 0.5 x
  \addplot[engblue, very thick] { 7.82 - 0.5*x };
  % marker at 1 GHz (x=9): delta = 2.1 um, log10 = 0.32
  \draw[engred, dashed] (axis cs:9,0) -- (axis cs:9,0.32);
  \node[engred, font=\scriptsize, anchor=south west] at (axis cs:9,0.4)
    {$1\,$GHz: $2.1\,\mu$m};
  % marker at 1 MHz (x=6): delta = 66 um, log10 = 1.82
  \draw[enggray, dashed] (axis cs:6,0) -- (axis cs:6,1.82);
  \node[enggray, font=\scriptsize, anchor=west] at (axis cs:6,2.2)
    {$1\,$MHz: $66\,\mu$m};
\end{axis}
\end{tikzpicture}
\caption[Skin depth of copper versus frequency]{The skin depth of copper,
$\delta = 1/\sqrt{\pi f \mu \sigma}$, on logarithmic axes, where the
inverse-square-root law plots as a straight line of slope $-\tfrac{1}{2}$.
The depth falls from tens of micrometres at megahertz frequencies to a few
micrometres at gigahertz frequencies, so a plating or shield of a few skin
depths suffices to carry RF current or to attenuate a field. At
$1\,\si{\giga\hertz}$ the depth is $2.1\,\si{\micro\meter}$, which is why
silver or gold plate only microns thick serves at microwave frequencies and
why solid-conductor thickness past a few skin depths buys nothing.}
\label{fig:vol04ch10:skin-depth}
\end{figure}
